\chapter{Συζήτηση και Συμπεράσματα}
\label{chapter:conc}

Σε αυτό το κεφάλαιο θα καταγραφούν οι τελικές σκέψεις πάνω στην παρούσα εργασία. Μέσα από εκτενή μελέτη και αναλυτικά πειράματα, έχω καταλήξει σε συμπεράσματα που αφορούν τόσο την πειραματική διαδικασία, όσο και το γενικό στόχο της εργασίας. Εκτός από αυτά τα πορίσματα, θα παρουσιαστούν και περιορισμοί, που συναντήθηκαν από εμένα συγκεκριμένα αλλά και του θέματος που πραγματεύεται η εργασία

\section{Συζήτηση}
Μέσω των πειράματα που πραγματοποιήθηκαν στο πλαίσιο της εργασίας (Κεφάλαιο 9), καταλήγουμε σε ορισμένα χρήσιμα συμπεράσματα για την απόδοση του νεοσύστατου μοντέλου. Η γενική εικόνα είναι ότι οι κανόνες βελτιώνουν το μοντέλο, αλλά το πλήρες σύνολο δεν φαίνεται το βέλτιστο και το όφελος εξαρτάται από το ρεπερτόριο. 

Οι κανόνες βελτιώνουν το μοντέλο στη ρίζα και στη συμπεριφορά. Το πλήρες μοντέλο υπερτερεί του baseline κατά 2,29 ποσοστιαίες μονάδες στη Joint ακρίβεια (από 63,15\% σε 65,43\%, [+1,62, +2,96]). Το κέρδος προέρχεται κυρίως από τη ρίζα (+1,96) και λιγότερο από την ποιότητα (+0,46). Ταυτόχρονα το σφάλμα μειώνεται και στις εννέα μετρικές της βιβλιογραφίας, έως 74\% στο CBS, και οι αποκλίσεις των κανόνων 1, 2, 3, 5 και 6 από τα δεδομένα πέφτουν κατά 79\% έως 96\%. Η βελτίωση της Joint ακρίβειας ξεχωρίζει από την τυχαιότητα της εκπαίδευσης, που μετρήθηκε σε τρία μοντέλα με δύο εκπαιδεύσεις το καθένα (τυπική απόκλιση 0,21 μονάδες), και επιβιώνει της διόρθωσης για πολλαπλές συγκρίσεις (Holm, p = 0,015).

Το baseline ξεφεύγει από τα δεδομένα στην ποικιλία και στον ρυθμό. Χρησιμοποιεί κατά μέσο όρο περίπου τρεισήμισι διαφορετικές συγχορδίες περισσότερες ανά τραγούδι από τα δεδομένα (bias του CC +3,53) και αλλάζει συγχορδία πολύ συχνά, στο 31,4\% των διαδοχικών νοτών έναντι 22,7\% στο full, και σε περισσότερες διαφορετικές θέσεις του μέτρου (bias του HRC +1,61). Οι κανόνες μειώνουν τα δύο bias στο +1,26 και στο +0,51, χωρίς να τα εξαλείφουν.

Η τήρηση ενός κανόνα και η ακρίβεια είναι διαφορετικοί στόχοι. Το μοντέλο με την υψηλότερη ακρίβεια, το full - {1,2,5} (66,4\%), τηρεί τον Κανόνα 1 λιγότερο από το full (απόκλιση 2,17 έναντι 0,60), αλλά έχει ισοδύναμη πιστότητα στις μετρικές (CBS 0,058 έναντι 0,059). Αντίστροφα, ο Κανόνας 6 πετυχαίνει τον στόχο του, αφού η απόκλισή του πέφτει από 11,21 σε 2,06 στο only {6,8}, αλλά η ομάδα Γ δίνει το μικρότερο κέρδος ακρίβειας (+1,24). Η ομοιότητα με τα στατιστικά των δεδομένων δεν μεταφράζεται λοιπόν πάντα σε συγχορδίες που ταιριάζουν με την αναφορά.

Τα κέρδη των ομάδων δεν προστίθενται. Μόνες τους, οι ομάδες Α, Β και Γ δίνουν +2,05, +2,97 και +1,24, δηλαδή +6,27 συνολικά, ενώ όλες μαζί δίνουν +2,29. Τα βάρη των κανόνων υπολογίστηκαν ξεχωριστά για τον καθένα, και στο πλήρες μοντέλο το άθροισμά τους (2,31 + 2,09 + 1,54 + 2,19 + 1,22 = 9,35) είναι αρκετές φορές μεγαλύτερο από το βάρος των δύο όρων cross-entropy. Μια πιθανή εξήγηση είναι ότι οι ομάδες διορθώνουν σε μεγάλο βαθμό το ίδιο έλλειμμα, την υπερβολική ποικιλία και τις πολύ συχνές αλλαγές του baseline, και ότι οι όροι ανταγωνίζονται μεταξύ τους και με τον κύριο στόχο. Συμφωνεί με αυτήν το ότι οι ομάδες Α και Β, μόνες τους, φέρνουν ήδη τις μετρικές ποικιλίας κοντά στις τιμές του full. Η εξήγηση δεν ελέγχθηκε όμως με πείραμα, για παράδειγμα με μικρότερα βάρη.

Η ομάδα Α κοστίζει στο πλήρες μοντέλο. Χωρίς τους κανόνες 1, 2 και 5 η ακρίβεια ανεβαίνει κατά 0,92 μονάδες [+0,25, +1,59] (p = 0,022, που όμως δεν επιβιώνει της διόρθωσης για πολλαπλές συγκρίσεις), και στο POP909 κατά περίπου 3,5 (32,68\% έναντι 29,20\%). Στο ίδιο σύνολο το only {1,2,5} έχει 27,28\%, κάτω από το baseline (27,68\%). Η ομάδα ζητά οι συγχορδίες να περιέχουν τη νότα της μελωδίας, με επιείκεια ανά τύπο χρόνου. Στην ποπ και στην τζαζ όμως οι νότες της μελωδίας είναι συχνά επεκτάσεις των συγχορδιών (ένατες, ενδέκατες) που ο πίνακας συμβατότητας δεν περιλαμβάνει για τις περισσότερες ποιότητες. Τα στατιστικά της ομάδας προέρχονται κυρίως από κλασικά τραγούδια, άρα είναι πιθανότατα λιγότερο κατάλληλα εκεί. Πρόκειται για υπόθεση που συμφωνεί με τα δεδομένα αλλά δεν ελέγχθηκε. 

Η ομάδα Β είναι η πιο αποτελεσματική μόνη της. Το only {3,7} ξεπερνά το πλήρες μοντέλο κατά 0,68 μονάδες, διαφορά οριακή (p = 0,077). Ο Κανόνας 7 επηρεάζει μία μόνο νότα ανά τραγούδι, άρα το κέρδος οφείλεται πιθανότατα στον Κανόνα 3, τη σταθερότητα του αρμονικού ρυθμού. Η ομαδοποίηση δεν επιτρέπει να το επιβεβαιώσουμε.

Οι συμπεριφορές που μετράμε δεν είναι ανεξάρτητες. Το only {1,2,5} μειώνει την απόκλιση του Κανόνα 3 από 11,35 σε 2,75 και του Κανόνα 6 από 11,21 σε 2,04, παρότι οι κανόνες 3 και 6 δεν είναι ενεργοί. Αντίστροφα, η απόκλιση του Κανόνα 3 σχεδόν δεν αλλάζει όταν αφαιρούνται οι κανόνες 3 και 7 από το full (1,88 έναντι 2,28). Ο Κανόνας 5 αφορά λίγες νότες και δεν μπορεί να κριθεί με αυτά τα δεδομένα. Η απόκλιση ενός κανόνα δεν αποδίδεται λοιπόν εύκολα στον ίδιο τον κανόνα.

Ο Κανόνας 7 υπερβάλλει. Τα μοντέλα με τον κανόνα τελειώνουν σε τονική στο 99,6\% έως 99,9\% των τραγουδιών, ενώ στα δεδομένα το ποσοστό είναι 73,7\%. Η ποινή 1 - P(τονική) ελαχιστοποιείται μόνο όταν η πιθανότητα είναι 1, και υπολογίζεται σε μία νότα ανά τραγούδι, όπου η cross-entropy έχει πολύ μικρό βάρος, οπότε ο όρος κυριαρχεί. Η υπερβολή είναι μεγαλύτερη όπου τα πραγματικά τραγούδια τελειώνουν συχνότερα αλλού: η απόκλιση του full είναι 24 μονάδες στο βασικό corpus, 49 στο CMD και 56 στο POP909. Μια διατύπωση που συγκρίνει την πιθανότητα με το ποσοστό των δεδομένων, όπως ο Κανόνας 3, ή που εξαρτάται από την πηγή, θα ταίριαζε καλύτερα.

Ο Κανόνας 8 προσφέρει λίγα. Η απόκλισή του είναι 3,42 στο baseline και 2,63 στο full, και στα υπόλοιπα μοντέλα κυμαίνεται από 2,26 έως 3,23, χωρίς ξεκάθαρη διάκριση. Συμπεραίνουμε λοιπόν, ότι η κατανομή της ποιότητας ανά λειτουργία μαθαίνεται σχεδόν εξίσου καλά από την cross-entropy μόνη της.

Η βαθμονόμηση σε κλασικά τραγούδια πιθανώς μεταφέρεται λιγότερο στην τζαζ και στην ποπ. Σε όλους τους κανόνες οι αποκλίσεις είναι μεγαλύτερες στο CMD και στο POP909 απ' ό,τι στο βασικό corpus (για τον Κανόνα 1 στο full: 0,57, 5,18 και 1,46). Οι κανόνες βαθμονομήθηκαν σε corpus όπου κυριαρχούν τα κλασικά τραγούδια. Το όφελος στην ακρίβεια είναι σαφές στο βασικό corpus (Joint από 69,03\% σε 71,48\%) και αβέβαιο στο CMD και στο POP909, όπου τα δείγματα επικύρωσης είναι μικρά (101 και 180 τραγούδια). Στις μετρικές όμως οι κανόνες μειώνουν το σφάλμα και εκεί: το CC κατά 50\% στο βασικό corpus, 34\% στο CMD και 47\% στο POP909. Η ποικιλία συγχορδιών της τζαζ και της ποπ μένει πιο δύσκολο πρόβλημα, αφού το MAE του CC στο full είναι 1,57, 5,77 και 7,28 αντίστοιχα.

Συγκριτικά με το paper αναφοράς, τα σχετικά σφάλματα του full είναι της ίδιας τάξης μεγέθους με του καλύτερου μοντέλου του paper και στις περισσότερες μετρικές μικρότερα, ενώ το baseline υστερεί στις CC, CTD, HRC και CBS. Στα τραγούδια του CMD το full είναι καλύτερο από το σύνολο τζαζ του paper σε πέντε μετρικές, ίσο περίπου σε τρεις και χειρότερο στη CC (48,4\% έναντι 31,5\%). Η σύγκριση όμως δεν είναι ισοδύναμη: το paper δίνει το καλύτερο από 16 διαμορφώσεις ανά μετρική, και τα δεδομένα, το πλέγμα, το λεξιλόγιο και η παραγωγή διαφέρουν. Δεν συμπεραίνουμε λοιπόν ότι το μοντέλο ξεπερνά το paper.

\section{Περιορισμοί}
Η έρευνα πάνω στο MIR και πιο συγκεκριμένα στην αυτόματη εναρμόνιση μελωδίας, αν και καινοτόμα και ενδιαφέρουσα, αποδείχτηκε πολύ απαιτητική. Οι περιορισμοί που συναντήθηκαν στην πορεία, ανάγονται και σε προσωπικό και ακαδημαικό επίπεδο, αλλά και στις δυσκολίες του επιλεγμένου θέματος.

\subsection{Περιορισμοί της εργασίας}
Οι παρακάτω περιορισμοί συνοψίζονται σε δύο κυρίως θέματα: το περιορισμένο χρονικό διάστημα σε συνδυασμό με τους διαθέσιμους πόρους. Σε ένα επαγγελματικό πλαίσιο με παραπάνω ελευθερία χρόνου και με παροχή υψηλότερων διαγραφών πόρους, θα μπορούσαν να πραγματοποιηθούν περισσότερα πειράματα και περισσότερη επεξεργασία των δεδομένων. Στο πλαίσιο όμως της διπλωματικής εργασίας, συναντήθηκαν οι παρακάτω περιορισμοί:

\begin{enumerate}
    \item Λίγες επαναλήψεις της εκπαίδευσης. Μόνο τρία μοντέλα (baseline, full και full - {1,2,5}) εκπαιδεύτηκαν δύο φορές, και τα υπόλοιπα πέντε μία. Η τυπική απόκλιση της εκπαίδευσης (0,21 μονάδες) εκτιμάται από τρία ζεύγη, με τρεις βαθμούς ελευθερίας και με την υπόθεση ότι είναι ίδια σε όλες τις ρυθμίσεις, άρα είναι αβέβαιη. Διαφορές μικρότερες από περίπου 0,6 μονάδες (τα full - {3,7}, full - {6,8} και only {1,2,5} έναντι του full) είναι εντός θορύβου.
    \item Τα διαστήματα εμπιστοσύνης του bootstrap υποτιμούν την αβεβαιότητα. Αφορούν μόνο την επιλογή των τραγουδιών επικύρωσης και όχι την τυχαιότητα της εκπαίδευσης, που καλύπτεται μόνο από τους ελέγχους t των τριών μοντέλων με δύο εκπαιδεύσεις.
    \item Πολλαπλές συγκρίσεις. Οι 11 συγκρίσεις ανάμεσα στα μοντέλα δεν είχαν προγραμματιστεί όλες εκ των προτέρων, και μετά τη διόρθωση Holm επιβιώνουν μόνο οι συγκρίσεις με το baseline. Το εύρημα ότι το full - {1,2,5} είναι καλύτερο από το full (p = 0,022 πριν από τη διόρθωση και 0,13 μετά) είναι λοιπόν ενδεικτικό.
    \item Μικρά δείγματα εκτός του βασικού corpus. Το CMD έχει 101 και το POP909 180 τραγούδια επικύρωσης, από τα 4.280 συνολικά, και οι διαφορές ανά πηγή είναι αβέβαιες.
    \item Δεν υπάρχει ανεξάρτητο σύνολο δοκιμής. Τα αποτελέσματα μετρήθηκαν στο ίδιο σύνολο επικύρωσης που επηρεάζει την εκπαίδευση, αφού η απώλεια επικύρωσης ρυθμίζει τον ρυθμό μάθησης. Αυτό έγινε γιατί τα δεδομένα συμβολικής μουσικής σε μία συγκεκριμένη μορφή, είναι δυσεύρετα. Ένας δεύτερος κίνδυνος θα ήταν να επιλεγεί η εποχή με το καλύτερο αποτέλεσμα επικύρωσης, γιατί τότε η ακρίβεια θα ήταν υπεραισιόδοξη. Αυτό δεν συμβαίνει: όλα τα μοντέλα εκπαιδεύονται για 200 εποχές και αξιολογείται πάντα η τελευταία.
    \item Η ομαδοποίηση των κανόνων. Για να περιοριστεί ο αριθμός των εκπαιδεύσεων, οι κανόνες εξετάστηκαν σε ομάδες ({1,2,5}, {3,7}, {6,8}), οπότε οι επιδράσεις δεν αποδίδονται σε μεμονωμένους κανόνες. Για να εκπαιδευτεί το μοντέλο σε όλους τους πιθανούς συνδυασμούς κανόνων, θα χρειαζόταν είτε πολύς παραπάνω χρόνος, είτε πολύ ακριβοί πόροι.
    \item Η βαθμονόμηση των κανόνων. Το βάρος κάθε κανόνα προκύπτει από ένα στατιστικό των δεδομένων που αφορά μόνο εκείνον, και δεν δοκιμάστηκε ποτέ ο συνδυασμός των βαρών ώστε να ρυθμιστούν μαζί. Αυτό για παράδειγμα μπορεί να βοηθούσε ώστε η διατύπωση του Κανόνα 7 να αναπαριστά πράγματι το ποσοστό των δεδομένων.
    \item Η «τήρηση» είναι ομοιότητα με τα δεδομένα. Αυτό σαν λειτουργία, βασίζεται στην πεποίθηση ότι τα δεδομένα τηρούν σε επιθυμητά επίπεδα τους κανόνες. Μικρή απόκλιση σημαίνει ότι το μοντέλο μιμείται τα δεδομένα και όχι ότι δεν παραβιάζει τον κανόνα.
    \item Όπως έχει αναφερθεί ήδη, η σύγκριση με το paper αναφοράς δεν είναι ισοδύναμη. Το paper δίνει το καλύτερο από 16 διαμορφώσεις ανά μετρική, ενώ εδώ υπάρχουν το πολύ δύο εκπαιδεύσεις ανά ρύθμιση. Διαφέρουν επίσης τα δεδομένα (το δικό μας σύνολο επικύρωσης προέρχεται από τις ίδιες πηγές με την εκπαίδευση, ενώ το σύνολο τζαζ του paper είναι εκτός εκπαίδευσης), το πλέγμα και το λεξιλόγιο συγχορδιών (13 × 24 έναντι 12 × 29) και η παραγωγή (μία διέλευση έναντι επαναληπτικής αποκάλυψης με δειγματοληψία).
    \item Τα δεδομένα με ευθυγραμμισμένες μελωδίες και συγχορδίες είναι λίγα και άνισα. Το βασικό corpus αποτελεί το 93,6\% των τραγουδιών, η τζαζ αντιπροσωπεύεται μόνο από τα 473 τραγούδια του CMD και η ποπ μόνο από το POP909, δηλαδή από κινεζικά ποπ τραγούδια. Η μείζονα τριφωνία καλύπτει το 56\% των νοτών, και αρκετές ποιότητες έχουν μόνο μερικές εκατοντάδες παραδείγματα. Τα βάρη κλάσεων μετριάζουν αλλά δεν εξαλείφουν το πρόβλημα.
\end{enumerate}

\subsection{Περιορισμοί του θέματος}
Οι παρακάτω περιορισμοί δεν οφείλονται στην υλοποίηση. Ισχύουν για κάθε προσέγγιση αυτόματης εναρμόνισης που αξιολογείται με παρόμοιο τρόπο.
 \begin{enumerate}
     \item Τα διαθέσιμα δεδομένα είναι λιγοστά. Ο χώρος του MIR και της επεξεργασίας συμβολικής μουσικής, δεν αφορά ένα πολύ ευρή κοινό. Οι έρευνες που έχουνε πραγματοποιηθεί πάνω στο αντικείμενο, δεν είναι σε πλήθος όσες υπάρχουν σε πιο δημοφιλή πεδία. Αυτό επηρεάζει και την ποικιλία και ποικιλομορφία των διαθέσιμων δεδομένων. Άλλος ένας παράγοντας που εμπλέκεται, είναι τα πνευματικά δικαιώματα των καλλιτεχνών, τα οποία αποτρέπουν πολλά μουσικά κομμάτια από το να ανεβαίνουν σε συμβολική μορφή δωρεάν στο διαδίκτυο.
     \item Οι μορφές αναπαράστασης της συμβολικής μουσικής είναι πολλές και διαφορετικές. Τα αρχεία μουσικών κομματιών είναι κωδικοποιημένα σε πάρα πολλά ξεχωριστά formats όπως: xml, musixml, midi, json και jams. Όλα αυτά διαφέρουν πολύ μεταξύ τους. Οπότε αν θέλει κανείς να χρησιμοποιήσει μία μεγάλη ενιαία βιβλιοθήκη συμβολικής μουσικής, θα πρέπει είτε να μπορέσει να τα μεταφέρει όλα επιτυχημένα σε ένα μοναδικό format, ή να κάνει το πρόγραμμά του προσαρμόσιμο σε όλα, το οποίο αυξάνει σε μεγάλο βαθμό τη δυσκολία.
     \item Δεν υπάρχει μία σωστή εναρμόνιση. Μια μελωδία δέχεται πολλές μουσικά αποδεκτές αρμονίες, αλλά τα δεδομένα περιέχουν μία. Η ακρίβεια ως προς τη συγχορδία αναφοράς θεωρεί λάθος κάθε εναλλακτική, οπότε η Joint ακρίβεια 65,4\% δεν σημαίνει ότι το 34,6\% των συγχορδιών είναι μουσικά λανθασμένο.
     \item Δεν υπάρχει καθιερωμένο μέτρο μουσικής ποιότητας. Οι εννέα μετρικές της βιβλιογραφίας και η τήρηση των κανόνων μετρούν ομοιότητα με στατιστικά πραγματικών τραγουδιών, όχι το πόσο καλά ακούγεται μια εναρμόνιση. Η υποκειμενική ποιότητα απαιτεί αξιολόγηση από ακροατές και μουσικούς.
     \item Οι κανόνες της θεωρίας είναι τάσεις με εξαιρέσεις και εξαρτώνται από το ύφος. Γι' αυτό εισάγονται ως μαλακοί περιορισμοί και όχι ως λογικοί. 
 \end{enumerate}

\section{Συμπεράσματα}
Η παρούσα εργασία, ξεκίνησε με σκοπό να μελετήσει αν μία νευροσυμβολική προσέγγιση στο πρόβλημα τις αυτόματης εναρμόνισης, θα συνδίαζε πράγματι τα πλεονεκτήματα των δύο αυτών κόσμων. Αν θα μπορούσε δηλαδή να κατακτήσει την ευφυΐα, την αναγνώριση μοτίβων και την εκφραστικότητα που διέπουν την μουσική σύνθεση, μέσω της τεχνολογίας των transformer, αλλά και την θεωρητικές βάσης που ακολουθεί κάθε κομμάτι μέσω της ενσομάτωσης των κανόνων. Τα συμπεράσματα ήταν ξεκάθαρα. Οι κανόνες μουσικής θεωρίας ως όροι της συνάρτησης απώλειας βελτιώνουν μετρήσιμα έναν Transformer εναρμόνισης μελωδίας, αλλά το όφελος εξαρτάται από τον κανόνα και από το ρεπερτόριο. Η εργασία κατασκεύασε ένα μοντέλο μόνο-encoder, περίπου 25,5 εκατομμυρίων παραμέτρων, που προβλέπει σε μία διέλευση τη ρίζα και την ποιότητα της συγχορδίας για κάθε νότα μιας μονοφωνικής μελωδίας. Στην εκπαίδευσή του οι κανόνες προστίθενται ως διαφορίσιμοι όροι ποινής, με βάρη που προκύπτουν από τα στατιστικά των δεδομένων, και το μοντέλο εκπαιδεύτηκε σε 21.504 τραγούδια από τα PDMX, OpenEWLD, CMD και POP909.

Όπως είδαμε και στα πειράματα, οι κανόνες πράγματι βελτιώνουν το μοντέλο. Η Joint ακρίβεια ανεβαίνει κατά 2,29 μονάδες, με διαφορά που ξεχωρίζει από την τυχαιότητα της εκπαίδευσης, και το σφάλμα μειώνεται και στις εννέα μετρικές της βιβλιογραφίας. Το πλήρες σύνολο κανόνων όμως, δεν φαίνεται το βέλτιστο. Τα κέρδη των ομάδων δεν προστίθενται (+6,27 μεμονωμένα έναντι +2,29 μαζί), και το μοντέλο που εκπαιδεύτηκε χωρίς την ομάδα {1,2,5} κατέγραψε υψηλότερη ακρίβεια από το πλήρες. Η διαφορά αυτή ξεχωρίζει από τον θόρυβο των εκπαιδεύσεων αλλά όχι από τη διόρθωση για πολλαπλές συγκρίσεις, ενώ οι υπόλοιπες ομάδες δεν ξεχώρισαν από το πλήρες. Αυτό μας φέρνει στο απλό συμπέρασμα ότι πρώτον, παραπάνω κανόνες δεν σημαίνει απαραίτητα καλύτερα αποτελέσματα. Για να φτιαχτεί ένα βέλτιστο μοντέλο εναρμόνισης μελωδίας στην συγκεκριμένη αρχιτεκτονική, θα πρέπει να πραγματοποιηθούν πολύ παραπάνω πειράματα με πολλαπλούς συνδυασμούς πολλών παραπάνω όρων. Αυτό ώστε όχι μόνο να βρεθούν μεμονωμένα οι κατάλληλοι κανόνες για αυτό το έργο, αλλά και και ο καλύτερος συνδυασμός τους. Επιπλέον, οι κανόνες λειτουργούν ως μαλακοί περιορισμοί με βάρη από τα δεδομένα και αποδίδουν όσο τα στατιστικά ταιριάζουν στο ρεπερτόριο. Το όφελος είναι σαφές στο βασικό corpus και αβέβαιο στα CMD και POP909, και η διατύπωση ενός όρου, όπως του Κανόνα 7, μετρά όσο και το βάρος του. Αυτό μας οδηγεί στο συμπέρασμα ότι οι κανόνες επηρεάζουν σε διαφορετικό βαθμό κάθε είδος μουσικής. Τα είδη με πιο περίπλοκες αρμονίες, όπως η Jazz και η Pop, φαίνεται ότι είτε χρειάζονται πιο εξειδικευμένους κανόνες πάνω στις ιδιαιτερότητές τους, είτε δεν επωφελούνται γενικά το ίδιο από τη χρήση κανόνων επειδή μπορεί να είναι πιο απρόβλεπτη η αρμονία τους. Μάλιστα οι κανόνες 6,8 που αφορούν την λειτουργική αρμονία, η οποία έχει ξεκινήσει από τη jazz, δεν ξεχωρίζουν στο CMD από τις άλλες ομάδες, που εμπνέονται περισσότερο από την κλασική αρμονία, αν και τα δείγματα είναι πολύ μικρά για να είναι βέβαιο.

\section{Μελλοντικές κατευθύνσεις}
Η παρούσα εργασία φανερά δίνει έναυσμα για πιθανές νέες μελέτες ή επεκτάσεις πάνω στο πρόβλημα το της εναρμόνισης μελοδίας. Με βάση λοιπόν όλα τα παραπάνω συμπεράσματα, καταλήγουμε στις εξής πιθανές μελοντικές κατευθήνσεις.

\begin{enumerate}
    \item Αρχικά το πιο προφανές είναι να εκπαιδευτούν περισσότερες φορές, με διαφορετικά seeds, τα πέντε μοντέλα που εκπαιδεύτηκαν μόνο μία. Μέσω περισσότερων επαναλήψεων κάθε ρύθμισης, μπορεί να μετρηθεί πιο αξιόπιστα η διακύμανση της εκπαίδευσης και να κριθούν οι οριακές διαφορές, ιδιαίτερα η διαφορά του full - {1,2,5} από το πλήρες μοντέλο.
    \item Παράλληλα, με αν υπήρχε περιθώριο για παραπάνω εκπαιδεύσεις, θα έπρεπε να ελεγχθούν οι κανόνες ένας ένας. Με Εκπαίδευση με μεμονωμένους κανόνες αντί για ομάδες, θα προέκυπταν αναλυτικότερα αποτελέσματα για το ποιοι από τους κανόνες επωφελούν σημαντικά το μοντέλο και ποιοι θα μπορούσαν να αντικατασταθούν.
    \item Σαν επόμενο βήμα, θα ήταν η εισαγωγή νέων κανόνων, εμπνευσμένοι ίσως από διαφορετικά είδη μουσικής. Θα έπρεπε προφανώς να ελεγχθούν με τον ίδιο τρόπο για το αν βοηθάν την απόδοση τη μοντέλου και αν μπορούν να συνδιαστούν με τους ήδη καθιερωμένους.
    \item Άλλη μία καινούρια προσέγγιση είναι η από κοινού ρύθμιση των βαρών των κανόνων. Θα ακολουθούσε δοκιμή με μικρότερα βάρη, για να ελεγχθεί η εξήγηση ότι τα κέρδη των ομάδων δεν προστίθενται επειδή οι όροι ανταγωνίζονται. Επίσης, με σταδιακή αύξηση των αυτών των βαρών, θα ελεγχόταν η οριακή τιμή τους στην οποία μεγιστοποιείται η συμβολή τους.
    \item Πολλά βασικά κενά της εργασίας θα κάλυπτε η συλλογή περισσότερων δεδομένων, ειδικλα για την τζαζ και την ποπ. Έτσι θα κριθεί αξιόπιστα το όφελος των κανόνων εκεί και σε επόμενο βήμα μπορεί να προστεθεί ανεξάρτητο σύνολο δοκιμής, κάτι βασικό που λείπει από την παρούσα μελέτη.
    \item Τέλος, για να μετρηθεί με τον ορθότερο τρόπο η υποκειμενική μουσική ποιότητα, που δεν καλύπτουν η ακρίβεια και οι μετρικές, πρέπει να ενταχθεί αξιολόγηση από ακροατές και πιο συγκεκριμένα μουσικούς. 
%Αναδιατύπωση του κανόνα 7
\end{enumerate}
